\documentclass[11pt]{article}

\usepackage[margin=1in]{geometry}
\usepackage{amsmath,amssymb,amsthm,mathtools}
\usepackage{booktabs,longtable,array}
\usepackage{graphicx}
\usepackage{xcolor}
\usepackage{hyperref}
\usepackage{listings}
\usepackage{microtype}

\hypersetup{
  colorlinks=true,
  linkcolor=blue!55!black,
  citecolor=blue!55!black,
  urlcolor=blue!55!black
}

\lstset{
  basicstyle=\ttfamily\small,
  columns=fullflexible,
  breaklines=true,
  frame=single,
  backgroundcolor=\color{black!3}
}

\newtheorem{theorem}{Theorem}
\newtheorem{proposition}[theorem]{Proposition}
\newtheorem{lemma}[theorem]{Lemma}
\newtheorem{corollary}[theorem]{Corollary}
\theoremstyle{definition}
\newtheorem{definition}[theorem]{Definition}
\newtheorem{hypothesis}[theorem]{Analytic input}
\theoremstyle{remark}
\newtheorem{remark}[theorem]{Remark}

\newcommand{\one}{\mathbf 1}
\newcommand{\RR}{\mathbb R}
\newcommand{\QQ}{\mathbb Q}
\newcommand{\TT}{\mathcal T}
\newcommand{\XiSet}{\boldsymbol\Xi}
\newcommand{\eps}{\varepsilon}
\newcommand{\etaP}{\eta}
\newcommand{\Hone}{H_1}
\newcommand{\hashsplit}[2]{\shortstack[l]{\ttfamily #1\\\ttfamily #2}}

\title{A Verification Dossier for a Computer-Assisted,\\
Manuscript-Conditional Bound \(H_1\leq 236\)}
\author{Research draft for independent verification}
\date{3 September 2026}

\begin{document}
\maketitle

\begin{abstract}
This document gives a verification-oriented proof of the following conditional
statement.  If Propositions 1--3 of Stadlmann's 2026 manuscript on bounded
prime gaps are correct, with the explicitly stated Bombieri--Vinogradov repair
of the negative-parameter branch in Proposition 3(D), then a finite
computer-assisted certificate gives
\[
  \Hone\leq 236.
\]
The finite witness has dimension \(k=48\), a two-slice count-capped support,
and three two-piece rational polynomial product profiles.  The six entries of
the \(I\)-matrix are exact rational numbers; the six entries of the
\(J\)-matrix are enclosed using 1024-bit Arb ball arithmetic.  The
minorant-adjusted lower endpoint is
\[
  1.001530068839069\ldots>1.
\]
All support inequalities are reduced either to explicit hand arguments or to
exact rational linear-programming dual checks.  The conclusion is not claimed
as an unconditional theorem: the analytic input is unrefereed and has not
been reproved here, and the new finite certificates still require independent
implementation-level reproduction.
\end{abstract}

\begin{center}
\fbox{\parbox{0.92\textwidth}{
\textbf{Status.}  The proof below is complete relative to Analytic Input
\ref{hyp:analytic} and the stated exact/interval computation model.  It must
not be cited as an unconditional improvement to the peer-reviewed bound
\(H_1\leq246\).  Section~\ref{sec:gaps} lists every presently known gap or
external dependency.
}}
\end{center}

\tableofcontents

\section{Statement and logical dependencies}

Let \(p_n\) be the \(n\)-th prime and put
\[
  H_1:=\liminf_{n\to\infty}(p_{n+1}-p_n).
\]
An integer set \(\mathcal H=\{h_1,\ldots,h_k\}\) is \emph{admissible} if,
for every prime \(p\), its reduction modulo \(p\) omits at least one residue
class.

The published Maynard sieve shows how a suitable variational inequality for a
symmetric function supported in an appropriate simplex implies that infinitely
many translates of an admissible tuple contain two primes
\cite{Maynard2015}.  Stadlmann's Proposition 1 is the version needed here: it
allows a union of support slices and a Harman-sieve minorant
\cite{Stadlmann2026}.  We isolate that analytic ingredient rather than silently
assuming the whole manuscript.

\begin{hypothesis}[Stadlmann sieve input, with a branch repair]
\label{hyp:analytic}
Propositions 1, 2 and 3 of \cite{Stadlmann2026} hold with the conventions used
in that manuscript.  In Proposition 3(D), the range \(\omega_0<0\) is handled
by the Bombieri--Vinogradov theorem, while the four-bin Type-IIc partition is
required only for \(0\leq\omega_0\leq\omega\).  The endpoint
\(\omega_0=0\) belongs to the latter range.
\end{hypothesis}

The repair is forced by the printed fourth capacity \(8\omega_0\): it is
negative when \(\omega_0<0\), so the printed partition condition cannot hold
literally there.  For negative \(\omega_0\), the relevant modulus exponent is
strictly below \(1/2\), which is the Bombieri--Vinogradov range.  A polished
version should incorporate this split into the statement of the analytic
proposition itself.

We use the following direct restatement of Stadlmann's main sieve implication.

\begin{proposition}[Imported sieve implication]
\label{prop:imported}
Let \(\rho(n;x)\) and constants \(c_1,c_2\) satisfy:
\begin{enumerate}
  \item \(-c_2\leq\rho(n;x)\leq\one_{\mathbb P}(n)\) on \([x,2x]\);
  \item \(\rho\) has the required equidistribution for the moduli generated by
        a support \(\TT_k(\delta,A,B,\eps)\);
  \item if \(\rho(n;x)\neq0\), every prime factor of \(n\) exceeds
        \(x^\beta\), with \(\beta>\max_j B_{j,1}\);
  \item \(\sum_{x\leq n\leq2x}\rho(n;x)=(1-c_1+o(1))x/\log x\).
\end{enumerate}
If a symmetric square-integrable \(F\), supported on
\(\TT_k(\delta,A,B,\eps)\), satisfies
\begin{equation}
\label{eq:sieve-ratio}
  \frac{k(1-c_1)J(F;\eps)-k c_2K(F;\eps)}{I(F;\eps)}>1,
\end{equation}
then every admissible \(k\)-tuple \(\mathcal H\) gives
\[
  H_1\leq \max\mathcal H-\min\mathcal H.
\]
\end{proposition}

The definitions of \(I,J,K\) are those of Definition 5 in
\cite{Stadlmann2026}.  In particular, \(I\) is the \(L^2\)-mass of \(F\),
\(J\) contains two last-coordinate integrations and \(k-1\) common
coordinates, and \(K\) is the penalty from the thin region where the common
coordinate sum lies past the relevant slice endpoint.

\begin{theorem}[Conditional computer-assisted result]
\label{thm:main}
Assume Analytic Input~\ref{hyp:analytic}.  Assume also that the exact rational
and Arb computations described in Sections~\ref{sec:support-cert} and
\ref{sec:integrals} are executed by conforming implementations of rational
arithmetic and outward-rounded ball arithmetic.  Then
\[
  \boxed{H_1\leq236}.
\]
\end{theorem}

The rest of the document proves Theorem~\ref{thm:main} and records enough
information to reproduce every finite computation.

\section{The two-slice support}
\label{sec:support}

All terminating decimals below denote exact rational numbers.  Put
\begin{equation}
\label{eq:parameters}
 k=48,\qquad \delta=0.014=\frac7{500},\qquad
 A_1=0.253=\frac{253}{1000},\qquad
 A_2=0.258=\frac{129}{500}.
\end{equation}
For \(t=(t_1,\ldots,t_{48})\in[0,1]^{48}\), define
\[
 S(t)=\sum_{i=1}^{48}t_i,\qquad
 m(t)=\#\{i:t_i>\delta\},\qquad
 L(t)=\sum_{t_i>\delta}t_i.
\]
The two cap sequences are
\begin{center}
\begin{tabular}{c|rrrrrrr}
\toprule
\(m\) & 0&1&2&3&4&5&\(6+\)\\
\midrule
\(B^{(1)}_m\) &0&.156&.156&.170&.184&.190&.202\\
\(B^{(2)}_m\) &0&.153&.153&.167&.181&.190&.190\\
\bottomrule
\end{tabular}
\end{center}
For \(m>6\), the last displayed entry is repeated.  The limiting support is
\begin{equation}
\label{eq:T0}
\begin{split}
 \TT_0={}&\{t:S(t)\leq A_1,\ L(t)\leq B^{(1)}_{m(t)}\}\\
 &\cup\{t:A_1<S(t)\leq A_2,\ L(t)\leq B^{(2)}_{m(t)}\}.
\end{split}
\end{equation}
Changing weak to strict inequalities on the boundary has no effect on any
integral.

For the positive enlargement parameter \(\eps\), use Stadlmann's Definition 1
with
\[
 A=(A_0,A_1,A_2)=(-\eps,0.253,0.258)
\]
and the two rows \(B^{(1)},B^{(2)}\).  Denote the resulting support by
\(\TT_\eps\).  Since \(B^{(1)}_m\geq B^{(2)}_m\) for every \(m\), one has
\begin{equation}
\label{eq:nesting}
 \TT_0\subseteq\TT_\eps\quad\text{up to boundary sets of measure zero.}
\end{equation}
Indeed, an upper-slice point moved into the first enlarged band encounters a
weaker, not a stronger, cap.

The cap conditions in Definition 1 hold:
\[
 \delta<B^{(j)}_m\leq B^{(j)}_{m+1}\leq B^{(j)}_m+\delta
 \qquad(m\geq1),
\]
with the evident convention for the plateaux.

\section{The minorant and global analytic inequalities}
\label{sec:analytic-checks}

Choose
\begin{equation}
\label{eq:xis}
 \xi_1=0.3795=\frac{759}{2000},\qquad
 \xi_2=\xi_3=0.4021=\frac{4021}{10000}.
\end{equation}
The five hypotheses of Stadlmann's Proposition 2 are checked exactly:
\begin{align*}
 2\xi_1+3\xi_2&=\frac{19653}{10000}<2,
   &2-(2\xi_1+3\xi_2)&=\frac{347}{10000},\\
 \xi_2&=\xi_3,\\
 \xi_1+9\xi_2&=\frac{2499}{625}<4,
   &4-(\xi_1+9\xi_2)&=\frac1{625},\\
 2\xi_1+\xi_2&=\frac{11611}{10000}>1,
   &(2\xi_1+\xi_2)-1&=\frac{1611}{10000},\\
 17\xi_2&=\frac{68357}{10000}<7,
   &7-17\xi_2&=\frac{1643}{10000}.
\end{align*}

Let \(\etaP=10^{-10}\) denote the fixed safety parameter printed in
Stadlmann's Proposition 3; this is distinct from the support enlargement
parameter \(\eps\).  Substitution of \eqref{eq:parameters}--\eqref{eq:xis}
into the global conditions of Proposition 3 gives the following exact slacks:
\begin{center}
\begin{tabular}{lll}
\toprule
Condition & Exact slack & Decimal value\\
\midrule
Type I, binding branch & \(2499997/15000000000\) & \(1.6666647\cdot10^{-4}\)\\
Type II range & \(3000001/100000000\) & \(0.03000001\)\\
Type II, first \(\delta\)-branch & \(3049999/5000000000\) & \(0.0006099998\)\\
Type II, second \(\delta\)-branch & \(124999/5000000000\) & \(0.0000249998\)\\
Type III & \(28187499/5000000000\) & \(0.0056374998\)\\
\bottomrule
\end{tabular}
\end{center}
Thus the tight global inequality is the second Type-II \(\delta\)-branch, but
it still has positive rational slack.

Proposition 2 produces a minorant with
\begin{equation}
\label{eq:beta-c2}
 \beta=1-2\xi_2=\frac{979}{5000}=0.1958,
 \qquad c_2=24.
\end{equation}
The largest one-large-coordinate cap is \(0.156\), and hence
\begin{equation}
\label{eq:beta-slack}
 \beta-\max_jB_{j,1}=\frac{199}{5000}=0.0398>0.
\end{equation}

\section{Certification of all support-cell partition conditions}
\label{sec:support-cert}

For cap bounds \(B_1,B_2\), counts \(m_1,m_2\), and threshold \(\delta\),
write
\[
 \XiSet(B_1,B_2,m_1,m_2,\delta)
 =\left\{y\in[\delta,1]^{m_1+m_2}:
 \sum_{i\leq m_1}y_i\leq B_1,
 \sum_{i>m_1}y_i\leq B_2\right\}.
\]
Stadlmann's Proposition 3 asks that every point of every such set admit the
partitions (A)--(E) printed there.

\subsection{Reduction to 105 frontier pairs}

If a cap sequence is constant from count six onward, any block of more than
six entries can be reduced to six entries by repeatedly replacing two entries
by their sum.  The merged entry remains in \([\delta,1]\), because the whole
block sum is at most \(0.202<1\).  A packing of the merged entries expands to a
packing of the original entries by putting the two constituents in the same
bin.  This reduces both cap rows to counts \(0,\ldots,6\).

There are therefore
\[
 \binom{7+1}{2}+7^2+\binom{7+1}{2}=28+49+28=105
\]
unordered lower/lower, lower/upper, and upper/upper cell pairs.

\subsection{Exact order-statistic witnesses}

For each block, sort its entries increasingly.  If \(r\) of the \(q\)
remaining entries are placed into a side bin, their sum is at most
\begin{equation}
\label{eq:order-bound}
 \frac{r}{q}\bigl(B-a\delta\bigr),
\end{equation}
where \(a\) already-removed entries contribute at least \(a\delta\).
For a fixed list of bin capacities, it is therefore enough to enumerate:
the primary bin, the order of the side bins, and the number of consecutive
smallest entries from each original block sent to each side bin.  Every bound
in \eqref{eq:order-bound} is rational.

The exact oracle in \texttt{primegaps.distribution} performs precisely this
finite enumeration.  On the final support it proves, for all 105 pairs,
conditions (A), (B), (C), and (E).  It also proves Type IIc condition (D) for
73 pairs.  Its remaining 32 failures are conservative failures of this fixed
order-statistic template, not counterexamples.

The executable check \texttt{analysis/verify\_final\_support.py} independently
enumerates the 105 pairs and asserts the exact split
\[
 105=(73\text{ direct Type-IIc})+(8\text{ hand})
       +(24\text{ exact-LP}).
\]
The final cap surgery makes the formerly reported lower/lower Type-IIb
exceptions disappear: the exact oracle now certifies all Type-IIb pairs
directly.  This corrects a stale statement in the exploratory report.

\subsection{The eight adaptive Type-IIc hand cases}
\label{sec:hand}

For Type IIc, put
\begin{align}
 C_1&=\gamma-2\delta-8\omega_0-\etaP,&
 C_2&=\tfrac12-\gamma-2\omega_0-\etaP,\nonumber\\
 C_3&=\delta+4\omega_0-\etaP,&
 C_4&=8\omega_0.\label{eq:Ccaps}
\end{align}
We work only with \(0\leq\omega_0\leq\omega\), as stipulated in Analytic
Input~\ref{hyp:analytic}.  Also
\begin{equation}
\label{eq:C2identity}
 C_2=0.472-10\omega_0-2\etaP-C_1.
\end{equation}

First consider an upper/upper pair with one block of count \(a\in\{1,2\}\)
and cap \(0.153\), and a second block of count \(m\in\{3,4,5\}\).
Keep the first block in bin 1.  Write the second block as
\(b_1\leq\cdots\leq b_m\), and let
\(r=\max(0,T-C_1)\) be the mass that must be removed from bin 1.  For
\(m=5\) we deliberately use the larger obsolete cap \(0.195\), so the proof
dominates the final cap \(0.190\).  Exact bounds are:
\begin{center}
\begin{tabular}{c|ccc}
\toprule
\(m\)&dominating \(T\)&upper bound for \(r\)&base term in \(C_2\)\\
\midrule
3&.320&\(0.0099000002\)&\(0.0719999998\)\\
4&.334&\(0.0239000002\)&\(0.0579999998\)\\
5&.348&\(0.0379000002\)&\(0.0439999998\)\\
\bottomrule
\end{tabular}
\end{center}
Here ``base term'' means \(C_2\geq\text{base}+r\), obtained from
\eqref{eq:C2identity}.

For \(m=3\), \(r<\delta\leq b_1\); move \(b_1\) to bin 2.  Since
\(b_1\leq0.167/3<0.056<C_2\), this works.

For \(m=4\), if \(b_1\geq r\), move \(b_1\) to bin 2; here
\(b_1\leq0.181/4<C_2\).  Otherwise
\(b_1+b_2\geq2\delta>r\).  Put \(b_2\) in bin 2 and \(b_1\) in bin 3.
The bounds
\[
 b_2\leq\frac{0.181-0.014}{3}<0.056<C_2
\]
hold.  Moreover \(r>b_1\geq\delta\) forces
\(\omega_0>0.0067624\), whence \(C_3>0.041\), while
\(b_1<r<0.024\).

For \(m=5\), remove \(b_1\) if it crosses \(r\).  Otherwise remove
\(b_1,b_2\) if they cross \(r\); put \(b_2\) in bin 2 and \(b_1\) in bin 4.
In this case \(r>\delta\) gives \(C_4>0.040\), whereas
\(b_1<r<0.038\).  In the remaining case
\(b_1+b_2<r<3\delta\), remove \(b_1,b_2,b_3\); put \(b_3\) in bin 2 and
\(b_1,b_2\) in bins 3 and 4.  Then
\[
 b_3\leq\frac{0.195-2\delta}{3}<0.056<C_2,
 \qquad b_2<r-\delta<0.024,
\]
while \(r>2\delta\) gives \(C_3>0.041\) and \(C_4>0.054\).
All margins are many orders of magnitude larger than \(\etaP\).

The two mixed hand cases are lower-3/upper-4 and lower-4/upper-3.  Their total
is at most \(0.351\), so
\[
 r<0.0209000002<2\delta,
 \qquad C_2\geq0.0659999998+r.
\]
Choose \(b_1\leq b_2\leq\cdots\) from the upper block.  If
\(b_1\geq r\), put it in bin 2.  In this case
\[
 b_1\leq\max(0.181/4,0.167/3)<0.056<C_2.
\]
Otherwise remove
\(b_1,b_2\), placing the larger in bin 2 and the smaller in bin 3.  Now
\[
 b_2\leq\max\left(\frac{0.181-\delta}{3},
                   \frac{0.167-\delta}{2}\right)
     =0.0765<C_2,
\]
because \(r>\delta\) makes \(C_2>0.080\).  The same inequality forces
\(\omega_0>0.0046374\), so \(C_3>0.0325>r>b_1\).

All inequalities in this subsection are replayed as exact fractions by
\texttt{analysis/verify\_adaptive\_}\allowbreak
\texttt{hand\_lemmas.py}.

\subsection{Exact rational dual certificates for the other 24 cases}
\label{sec:dual}

For a difficult pair, ignore its two-block origin and enlarge it to the set of
globally sorted masses
\[
 \delta\leq y_1\leq\cdots\leq y_n,
 \qquad \sum_{i=1}^ny_i\leq T.
\]
The value of \(T\) in Appendix~\ref{app:dual-jobs} is at least the sum of the
two actual caps, so proving a packing on this enlarged set is sufficient.

Subdivide the rational rectangle
\[
 0\leq\omega_0\leq\omega,
 \qquad \xi_2-\etaP\leq\gamma\leq
 \frac13+8\omega+\frac73\delta+3\etaP
\]
into \(N^2\) rational rectangles.  For fixed parameters, let \(j\) be the
first stopping rank for which
\[
 \sum_{i=j+1}^n y_i\leq C_1,
 \qquad
 \sum_{i=j}^n y_i\geq C_1.
\]
For a fixed rectangle and rank \(j\), these inequalities, the ordering
constraints, the total constraint, and the parameter bounds define a rational
polytope \(P\).

The tail \(y_{j+1},\ldots,y_n\) is left in bin 1.  A finite search proposes a
partition of \(y_1,\ldots,y_j\) among the four bins.  For each bin the worst
load minus its capacity is a linear functional on \(P\).  The replayer proves
that its maximum is nonpositive by checking an LP dual with exact
\texttt{Fraction} arithmetic.  Specifically it verifies:
\begin{enumerate}
  \item the sign of every dual multiplier;
  \item exact stationarity in every primal variable;
  \item the exact rational dual objective;
  \item the desired bin inequality.
\end{enumerate}
Floating-point HiGHS is used only to propose the partition and dual
multipliers.  No floating-point comparison accepts a certificate.  Empty
stopping cells are decided by an exact feasibility criterion.  Closed
parameter rectangles overlap on their boundaries, so their union covers the
whole Type-IIc parameter range.

The 24 required jobs are listed in Appendix~\ref{app:dual-jobs}.  Three older,
stronger mixed jobs are also retained in the archive.  Across all 27 saved
replays there are 11,359 nonempty conditional polytopes and 45,436 exact bin
inequalities.  Equality is allowed in the Type-IIc capacities, so a zero
weakest slack at a rational cell boundary is legitimate.

\begin{proposition}[Finite support certificate]
\label{prop:support}
Subject to the negative-\(\omega_0\) split in Analytic
Input~\ref{hyp:analytic}, every local partition hypothesis (A)--(E) of
Stadlmann's Proposition 3 holds for the support \(\TT_\eps\).
\end{proposition}

\begin{proof}
Counts reduce to the 105 frontier pairs.  Exact order-statistic witnesses prove
(A), (B), (C), and (E) for every pair and (D) for 73.  Section~\ref{sec:hand}
proves eight further instances of (D).  Section~\ref{sec:dual} proves the
remaining 24.  These three classes are disjoint and exhaust all 105 pairs.
\end{proof}

By Proposition~\ref{prop:support}, the global checks in
Section~\ref{sec:analytic-checks}, and Analytic Input~\ref{hyp:analytic}, the
minorant supplied by Stadlmann's Proposition 2 has the equidistribution
required in Proposition~\ref{prop:imported}.

\section{The explicit rational trial function}
\label{sec:trial}

Put
\[
 R=A_2=\frac{129}{500},\qquad D=\delta=\frac7{500},
 \qquad L=\frac{101}{500}=0.202.
\]
Let \(T_n\) denote the \(n\)-th Chebyshev polynomial of the first kind and set
\[
 x_1(t)=\frac{2t}{D}-1,
 \qquad
 x_2(t)=\frac{2t-D-L}{L-D}.
\]
For \(\ell=1,2,3\), define
\begin{equation}
\label{eq:gprofile}
 g_\ell(t)=
 \begin{cases}
   \displaystyle\sum_{n=0}^{8}a_{\ell,1,n}T_n(x_1(t)),&0\leq t\leq D,\\[4pt]
   \displaystyle\sum_{n=0}^{12}a_{\ell,2,n}T_n(x_2(t)),&D<t\leq L,\\[4pt]
   0,&\text{otherwise}.
 \end{cases}
\end{equation}
The exact terminating-decimal coefficients \(a_{\ell,r,n}\) appear in
Appendix~\ref{app:coefficients}.  They are not error bounds for an underlying
smooth function: equation \eqref{eq:gprofile} is the exact definition of the
proof witness.

Let
\begin{equation}
\label{eq:vector}
 c=\left(\frac{571087}{10^9},
          -\frac{12804621}{125000000},1\right),
\end{equation}
and define
\begin{equation}
\label{eq:F}
 F(t_1,\ldots,t_{48})=
 \one_{\TT_0}(t_1,\ldots,t_{48})
 \sum_{\ell=1}^3c_\ell\prod_{i=1}^{48}g_\ell(t_i).
\end{equation}
This function is bounded, symmetric, square-integrable, and supported on
\(\TT_\eps\) for every positive \(\eps\), by \eqref{eq:nesting}.  Possible
jumps at \(D\) and at support boundaries occur on null sets and are harmless.

\section{Exact and interval evaluation of the key integrals}
\label{sec:integrals}

Write \(I_{ab}\) and \(J_{ab}\) for the bilinear matrix entries associated to
the product profiles \(\prod_i g_a(t_i)\) and \(\prod_i g_b(t_i)\), including
the indicator of \(\TT_0\).  It is convenient to report normalized matrices
\[
 \widehat I_{ab}=R^{-48}I_{ab},
 \qquad
 \widehat J_{ab}=R^{-49}J_{ab}.
\]
Consequently
\begin{equation}
\label{eq:scaling}
 \frac{48J(F;0)}{I(F;0)}
 =48R\frac{c^{\!T}\widehat Jc}{c^{\!T}\widehat Ic}.
\end{equation}

\subsection{Exact evaluation of \(I\)}

For fixed \(a,b\), let \(h(t)=g_a(t)g_b(t)\), with polynomial pieces
\(h_s\) on \([0,D]\) and \(h_b\) on \([D,L]\).  If
\(h_s(t)=\sum_ec_et^e\), then
\begin{equation}
\label{eq:small-laplace}
 \int_0^D h_s(t)e^{-st}\,dt
 =\sum_e\frac{c_e e!}{s^{e+1}}
 \left(1-e^{-Ds}\sum_{q=0}^e\frac{(Ds)^q}{q!}\right).
\end{equation}
Thus the transform is a polynomial in \(s^{-1}\) and \(E=e^{-Ds}\).
After taking the \((48-m)\)-th power, the coefficient of \(E^j\) represents a
shift by \(jD\).  For a large coordinate, expand
\(h_b(D+y)=\sum_ed_ey^e\); then
\begin{equation}
\label{eq:large-laplace}
 \int_D^\infty h_b(t)e^{-st}\,dt
 =e^{-Ds}\sum_e\frac{d_e e!}{s^{e+1}}.
\end{equation}
Extending the polynomial to infinity in \eqref{eq:large-laplace} is harmless:
the large-coordinate-sum cap is at most \(L\), so no accepted point has an
individual large coordinate exceeding \(L\).

The inverse transform rules
\[
 \sum_{p\geq0}\frac{a_p}{s^p}
 \longleftrightarrow
 \sum_{p\geq0}\frac{a_p y^p}{p!},
 \qquad
 \sum_{q\geq1}\frac{b_q}{s^q}
 \longleftrightarrow
 \sum_{q\geq1}\frac{b_q x^{q-1}}{(q-1)!}
\]
give respectively a small-sum CDF and a large-sum density.  Every status
\(m\), shift \(j\), total cap \(H\), and large cap \(M\) then reduces to
\begin{equation}
\label{eq:capped-convolution}
 \int_0^{\min(H,M)}D_m(x)A_{48-m,j}(H-x)\,dx,
\end{equation}
an exact rational polynomial integral.

The union \eqref{eq:T0} is evaluated without overlap by
\begin{align*}
 &\one_{S\leq A_1,L\leq B^{(1)}_m}
 +\one_{S\leq A_2,L\leq B^{(2)}_m}
 -\one_{S\leq A_1,L\leq B^{(2)}_m}.
\end{align*}
Summing \eqref{eq:capped-convolution} with multiplicity \(\binom{48}{m}\)
therefore gives every \(I_{ab}\) exactly.  FLINT \texttt{fmpq} arithmetic
produces
\begin{equation}
\label{eq:Imatrix}
\widehat I=\begin{pmatrix}
1.1343249249219832\!\cdot10^{-104}&2.0748726190460893\!\cdot10^{-106}&2.3278789274966597\!\cdot10^{-108}\\
2.0748726190460893\!\cdot10^{-106}&5.0100827747025920\!\cdot10^{-108}&8.2666312941118770\!\cdot10^{-110}\\
2.3278789274966597\!\cdot10^{-108}&8.2666312941118770\!\cdot10^{-110}&2.2300813056765934\!\cdot10^{-111}
\end{pmatrix}.
\end{equation}
The displayed decimals are only summaries.  The six authoritative values are
the exact numerator/denominator pairs in
\texttt{exact\_chebyshev\_i\_8\_12.json}.

\subsection{Rigorous evaluation of \(J\)}

For \(J_{ab}\), 47 common coordinates remain.  Condition on the number \(m\)
of common coordinates exceeding \(D\), and compress them to two residual
sums:
\[
 x=\text{large residual sum},\qquad
 z=\text{small residual sum}.
\]
For each profile, the last coordinate has four possible components: two total
bands, and small or large last-coordinate status.  Its exact marginal is a
piecewise polynomial in \((x,z)\).  Every endpoint comparison is a rational
horizontal, vertical, or diagonal line.  Taking the global arrangement of all
four components before multiplying the two marginals gives 2,655 rational
cells over all statuses, with at most 37 cells in a single status.

On each polygonal cell, the integrand is
\[
 D_m(x)D_{47-m}(z)P_{ab}(x,z),
\]
where all coefficients and polygon vertices are rational.  Vertical slicing
expresses both polygon edges as rational affine functions of \(x\); integrating
first in \(z\) and then in \(x\) reduces the cell to polynomial evaluation.
The geometry and branching are exact.  The polynomial contractions use
1024-bit Arb balls with outward rounding, because fully reduced rational
composition creates prohibitive coefficient swell.  The resulting enclosures
are
\begin{equation}
\label{eq:Jmatrix}
\resizebox{\textwidth}{!}{$\displaystyle
\widehat J\in\begin{pmatrix}
[8.62740627938\!\cdot10^{-106}\pm4.86\!\cdot10^{-118}]&[1.65374676917\!\cdot10^{-107}\pm3.31\!\cdot10^{-119}]&[1.9406722244\!\cdot10^{-109}\pm2.63\!\cdot10^{-120}]\\
[1.65374676917\!\cdot10^{-107}\pm3.31\!\cdot10^{-119}]&[4.01956058425\!\cdot10^{-109}\pm7.33\!\cdot10^{-121}]&[6.4891996669\!\cdot10^{-111}\pm6.97\!\cdot10^{-122}]\\
[1.9406722244\!\cdot10^{-109}\pm2.63\!\cdot10^{-120}]&[6.4891996669\!\cdot10^{-111}\pm6.97\!\cdot10^{-122}]&[1.569136856\!\cdot10^{-112}\pm3.30\!\cdot10^{-122}]
\end{pmatrix}. $}
\end{equation}

Contracting \eqref{eq:Imatrix}--\eqref{eq:Jmatrix} against the exact vector
\eqref{eq:vector} gives
\begin{align}
 c^{\!T}\widehat Ic
 &=1.994852232725948450897686\ldots\cdot10^{-110}
   \quad\text{(exact rational)},\label{eq:Iquad}\\
 c^{\!T}\widehat Jc
 &\in[1.6134397849\cdot10^{-111}\pm8.12\cdot10^{-122}].\label{eq:Jquad}
\end{align}
Using \eqref{eq:scaling}, outward-rounded ball arithmetic gives
\begin{equation}
\label{eq:rawball}
 48R\frac{c^{\!T}\widehat Jc}{c^{\!T}\widehat Ic}
 \in[1.0016199680\pm8.30\cdot10^{-11}].
\end{equation}

\subsection{A rational upper bound for the minorant loss}
\label{sec:loss}

Let \(t=\xi_2=4021/10000\) and \(\beta=1-2t=979/5000\).  Stadlmann's
constant \(c_1=L_A+L_B\) is the sum of two four-dimensional integrals.  In
each integral the integrand is the reciprocal of five prime-exponent factors.
Every factor is at least \(\beta\): the four displayed integration variables
are at least \(\beta\), and the remaining exponent is constrained to exceed
the last displayed prime exponent.  Hence each integrand is at most
\(\beta^{-5}\).

For region \(A\), the defining inequalities include
\(\beta<\alpha_4<\alpha_3<\alpha_2<\alpha_1<t\) and
\(\alpha_1+\alpha_2<t\).  Since \(\alpha_2\geq\beta\), one has
\(\alpha_1\leq t-\beta=3t-1\); hence the \(\alpha_1\)-width above
\(\beta\) is at most \(5t-2\).  Also
\(2\alpha_2\leq\alpha_1+\alpha_2<t\), and the same bound propagates
to \(\alpha_3,\alpha_4\).  Thus the four interval widths are at most
\[
 5t-2=\frac{21}{2000},\qquad
 \frac52t-1=\frac{21}{4000}
 \quad\text{(three times)}.
\]
For region \(B\), all four displayed variables are bounded above by
\(8t-3\).  Moreover \(\alpha_2,\alpha_4>\alpha_3\) gives
\(2\alpha_3<\alpha_2+\alpha_4<t\), and
\(\alpha_2+\alpha_4<t\) with \(\alpha_2\geq\beta\) gives
\(\alpha_4<t-\beta=3t-1\).  The four widths are therefore at most
\[
 10t-4=\frac{21}{1000}\quad\text{(twice)},\qquad
 \frac52t-1=\frac{21}{4000},\qquad
 5t-2=\frac{21}{2000}.
\]
Therefore
\begin{align}
 L_A&\leq
 \frac{(5t-2)(\frac52t-1)^3}{\beta^5}
 =\frac{75969140625}{14389093733134384},\nonumber\\
 L_B&\leq
 \frac{(10t-4)^2(\frac52t-1)(5t-2)}{\beta^5}
 =\frac{75969140625}{899318358320899},\nonumber\\
 c_1&\leq\overline c_1:=
 \frac{1291475390625}{14389093733134384}
 <8.976\cdot10^{-5}.
\label{eq:c1bound}
\end{align}
This deliberately crude rectangle majorant is rigorous and is sufficient for
the final crossing.

\subsection{The \(K\)-term and the positive-\(\eps\) limit}
\label{sec:Klimit}

At \(\eps=0\), the region defining \(K\) is empty wherever its integrand is
nonzero.  Indeed, suppose a point where \(F\neq0\) belongs to a band ending at
\(A_j\).  Then
\[
 \sum_{i=1}^{47}t_i\leq A_j
\]
because \(t_{48}\geq0\).  The strict \(K\)-condition instead requires the
common sum to exceed the maximum of the relevant band endpoints.  Hence
\begin{equation}
\label{eq:Kzero}
 K(F;0)=0.
\end{equation}

Combining \eqref{eq:rawball}, \eqref{eq:c1bound}, and the lower endpoint only,
we obtain
\begin{equation}
\label{eq:adjusted}
 48R(1-\overline c_1)
 \frac{c^{\!T}\widehat Jc}{c^{\!T}\widehat Ic}
 \geq1.0015300688390691785\ldots>1.
\end{equation}

It remains to pass from the limiting support to a positive enlargement
parameter, as required by the analytic proposition.  Keep the bounded function
\(F\) in \eqref{eq:F} fixed.  By \eqref{eq:nesting}, its \(I\)-integral is
unchanged for sufficiently small positive \(\eps\).  The indicator regions in
\(J(F;\eps)\) converge pointwise almost everywhere to those in \(J(F;0)\), and
the integrands have a fixed bounded compactly supported majorant.  Dominated
convergence gives
\[
 J(F;\eps)\longrightarrow J(F;0).
\]
The \(K\)-region is contained in strips whose common-sum width tends to zero;
absolute continuity of the Lebesgue integral gives
\[
 K(F;\eps)\longrightarrow K(F;0)=0.
\]
The strict margin in \eqref{eq:adjusted} therefore implies that, for some
sufficiently small \(\eps>0\),
\[
 \frac{48(1-c_1)J(F;\eps)-48\cdot24K(F;\eps)}{I(F;\eps)}>1.
\]
This is exactly \eqref{eq:sieve-ratio}.

\section{The admissible 48-tuple and completion of the proof}
\label{sec:tuple}

Consider
\begin{align*}
\mathcal H=\{&0,6,8,14,18,24,26,48,50,54,56,60,66,68,74,78,80,84,\\
&90,96,98,104,110,116,120,126,134,138,144,150,158,164,168,176,\\
&180,186,188,194,200,204,206,210,216,224,228,230,234,236\}.
\end{align*}
It contains 48 distinct offsets and has diameter 236.  For every prime
\(p\leq48\), the following residue is omitted:
\begin{center}
\begin{tabular}{c|rrrrrrrrrrrrrrr}
\toprule
\(p\)&2&3&5&7&11&13&17&19&23&29&31&37&41&43&47\\
\midrule
omitted residue&1&1&2&2&9&10&4&13&13&15&1&7&10&2&5\\
\bottomrule
\end{tabular}
\end{center}
For \(p>48\), 48 offsets cannot cover all \(p\) residue classes.  Thus
\(\mathcal H\) is admissible.

\begin{proof}[Proof of Theorem~\ref{thm:main}]
Sections~\ref{sec:analytic-checks} and \ref{sec:support-cert}, together with
Analytic Input~\ref{hyp:analytic}, supply a minorant satisfying all hypotheses
of Proposition~\ref{prop:imported}.  Sections~\ref{sec:trial}--\ref{sec:Klimit}
supply a symmetric square-integrable trial function satisfying the strict
variational inequality for some \(\eps>0\).  Applying
Proposition~\ref{prop:imported} to the explicit admissible tuple above yields
\[
 H_1\leq\max\mathcal H-\min\mathcal H=236.
\]
\end{proof}

\section{Reproduction protocol}
\label{sec:reproduction}

The companion packet uses the following layout:
\begin{lstlisting}
Small_Gaps_236_Conditional_Proof.tex
Small_Gaps_Improvement_Investigation.md
analysis/
tmp/primegaps_exact/
\end{lstlisting}
The directory \texttt{tmp/primegaps\_exact} is the clean repository snapshot
at commit
\begin{center}
\texttt{eaf6a9dc9276ece44807bdce48ae639b462fbf13}.
\end{center}
It is included because \texttt{exact\_chebyshev\_j.py} imports its rational
polygon machinery.  The earlier archive omitted this dependency.

From the root of the unpacked packet, run the inexpensive checks:
\begin{lstlisting}[language=bash]
python analysis/verify_final_support.py
python analysis/verify_adaptive_hand_lemmas.py
python analysis/check_chebyshev_certificate.py
python analysis/audit_bundle.py
\end{lstlisting}
Expected decisive lines include:
\begin{lstlisting}
PASS: 105 frontier pairs; all Type-I/IIa/IIb/III witnesses pass;
      Type-IIc splits as 73 oracle + 8 hand + 24 exact-LP
PASS: replay manifest covers 11,359 conditional polytopes and
      45,436 exact dual inequalities
PASS: rigorous lower endpoint exceeds one
PASS: explicit 48-tuple is admissible and has diameter 236
\end{lstlisting}

The quick audit validates saved outputs and hashes; it does not regenerate the
expensive computations.  To regenerate all exact Type-IIc dual checks, run
\begin{lstlisting}[language=bash]
python analysis/verify_final_support.py --full-dual-replay
\end{lstlisting}
To regenerate all six exact \(I\) entries and six Arb \(J\) entries, use:
\begin{lstlisting}[language=bash]
python analysis/exact_chebyshev_i.py \
  analysis/two_piece_chebyshev_8_12.json --all \
  --output /tmp/recomputed_I.json

for pair in 0,0 1,0 1,1 2,0 2,1 2,2; do
  compact=${pair/,/}
  python analysis/exact_chebyshev_j.py \
    analysis/two_piece_chebyshev_8_12.json \
    --pair "$pair" --arb-bits 1024 \
    --output "/tmp/recomputed_J${compact}.json"
done
\end{lstlisting}
The \(J\) computations are substantially more expensive than the final
contraction.  Python 3, NumPy, SciPy, and \texttt{python-flint} are required.

The authoritative variational input hashes are:
\begin{center}
\small
\begin{tabular}{lp{0.57\textwidth}}
\toprule
File & SHA-256\\
\midrule
\texttt{two\_piece\_chebyshev\_8\_12.json}&
\hashsplit{18955d99bb31582789f283be3fa112aaf}{37c308a8db9c362afc54beb1ef7f5aa}\\
\texttt{exact\_chebyshev\_i\_8\_12.json}&
\hashsplit{227890cf5ff73624b1d6bafdf87d71984}{d94b05469075f253e7f3f76f0b3b40c}\\
\texttt{arb\_chebyshev\_j00\_8\_12.json}&
\hashsplit{acabdaeb2370c0fd4da7e8cb8b60738e}{cc53e34ab7c8aa8916a473c187541be1}\\
\texttt{arb\_chebyshev\_j10\_8\_12.json}&
\hashsplit{55024779acd8fab474be227a02cb7d895}{8303ba124bcb99ca2e2efe61b961444}\\
\texttt{arb\_chebyshev\_j11\_8\_12.json}&
\hashsplit{e1b6f0bb8c935d08968b63aa2a1fc12}{1adf20b477baad7acbf1e2764bc30e3e7}\\
\texttt{arb\_chebyshev\_j20\_8\_12.json}&
\hashsplit{6ccf709bef7347354fc7592ad0c33ae0}{7605b2a270b4d56f2705b920c929ef27}\\
\texttt{arb\_chebyshev\_j21\_8\_12.json}&
\hashsplit{81c6cc56732d645ce36fa782fe0b2b7e}{af9d5020e1a0af8845b6de5dd8cc75e8}\\
\texttt{arb\_chebyshev\_j33\_8\_12.json}&
\hashsplit{ee4eee7d28c5dc4df652ed1a1e14f75}{8907c694a6615ac3a5e382ce79db0ef40}\\
\bottomrule
\end{tabular}
\end{center}

\section{Known gaps and items requiring further investigation}
\label{sec:gaps}

\begin{longtable}{p{0.04\textwidth}p{0.25\textwidth}p{0.22\textwidth}p{0.35\textwidth}}
\toprule
ID&Issue&Present status&Required next check\\
\midrule
\endfirsthead
\toprule
ID&Issue&Present status&Required next check\\
\midrule
\endhead
G1&Stadlmann's analytic framework is a v1 arXiv manuscript, not a refereed
theorem.&The finite calculation assumes Propositions 1--3.&Independently reprove
the Type I/II/III equidistribution estimates and the Harman-minorant passage,
including all uniformity and endpoint conventions.\\
\addlinespace
G2&Proposition 3(D) is literally impossible for negative \(\omega_0\) because
one capacity is \(8\omega_0<0\).&This document explicitly replaces that branch
by Bombieri--Vinogradov and keeps \(\omega_0=0\) in Type IIc.&Write this split
into a corrected formal proposition and verify the dyadic endpoint at
\(\omega_0=0\) with the manuscript's exact modulus conventions.\\
\addlinespace
G3&The support extension is computer-assisted rather than a conventional short
lemma.&Every accepted inequality is rational; 24 cases require exact LP duals.&
Export and archive the individual dual records, or replace them by a smaller
human-readable family of parametric packings.  The saved logs contain canonical
digests, not the full dual vectors.\\
\addlinespace
G4&The \(J\)-certificate relies on a nontrivial exact-geometry implementation.&
Rational geometry plus 1024-bit Arb gives tight balls; the pinned dependency is
now included.&Reimplement the six entries independently, preferably using a
Bernstein or interval-convolution representation.\\
\addlinespace
G5&The inexpensive audit reads saved \(I\) and \(J\) data and parses saved dual
logs.&It detects corruption and inconsistent parameters but does not repeat the
expensive construction.&Run the full dual replay and all six \(J\) jobs on a
second machine and compare hashes/enclosures.\\
\addlinespace
G6&The proof uses a limiting support at \(\eps=0\), while the sieve theorem uses
positive \(\eps\).&Section~\ref{sec:Klimit} gives the dominated-convergence and
absolute-continuity argument.&Check the exact open/closed endpoint conventions
against any revised manuscript version.\\
\addlinespace
G7&The rectangle bound for \(c_1\) is much larger than the numerical loss.&It is
still small enough for the crossing and is derived in Section~\ref{sec:loss}.&
Verify directly from the manuscript's two displayed exceptional regions that
the five denominator factors are all at least \(\beta\); this is simple but
belongs in an external line-by-line audit.\\
\addlinespace
G8&No independent formal proof assistant checks the Python/FLINT/Arb code.&The
proof uses standard exact and directed-rounding libraries.&Translate the final
rational polynomial identities and LP duals into Lean, Isabelle, or a minimal
standalone checker if theorem-level assurance is desired.\\
\bottomrule
\end{longtable}

The two highest-priority tasks are G1/G2 (the analytic theorem) and G3/G4 (an
independent finite-certificate implementation).  Until they are completed, the
correct label remains ``computer-assisted and manuscript-conditional.''

\appendix

\section{Exact Chebyshev coefficients}
\label{app:coefficients}

Every entry in the following tables is a terminating decimal interpreted as
an exact rational number.  The first table gives \(a_{\ell,1,n}\) on
\([0,D]\); the second gives \(a_{\ell,2,n}\) on \((D,L]\).

\begin{center}
\resizebox{\textwidth}{!}{%
\begin{tabular}{r|rrr}
\toprule
\(n\)&\(a_{1,1,n}\)&\(a_{2,1,n}\)&\(a_{3,1,n}\)\\
\midrule
0& 0.33994185099295593& 0.31439997663976893& 0.28933450386539988\\
1&-0.31473632235620519&-0.32797545323372218&-0.33717965695054575\\
2& 0.16080998699310434& 0.17105855001494782& 0.18127667696344230\\
3&-0.084756677283129689&-0.089197949353752748&-0.094460826889808960\\
4& 0.045321339945169521& 0.046475224429408317& 0.048382766928134827\\
5&-0.024363383571296741&-0.024144689815926913&-0.024472504025804882\\
6& 0.012998606098755659& 0.012408253795148046& 0.012184534035083159\\
7&-0.0066374594469853533&-0.0061162556640771181&-0.0058200857857910538\\
8& 0.0027935872469380140& 0.0025079207325502187& 0.0023295083474845499\\
\bottomrule
\end{tabular}}
\end{center}

\begin{center}
\resizebox{\textwidth}{!}{%
\begin{tabular}{r|rrr}
\toprule
\(n\)&\(a_{1,2,n}\)&\(a_{2,2,n}\)&\(a_{3,2,n}\)\\
\midrule
0& 0.043461710152066277& 0.027310154517352539& 0.016210093065552517\\
1&-0.040891450010470079&-0.030972829714833009&-0.021641016774680666\\
2& 0.021219827767331754& 0.017563329821976396& 0.013449285925462764\\
3&-0.011361288940722516&-0.009959337886991930&-0.0081233423464629091\\
4& 0.0061784459117329152& 0.005647376332586101& 0.0048307020741936661\\
5&-0.0033914853873386812&-0.0032021387179028901&-0.0028443818081275310\\
6& 0.0018731376364590881& 0.0018153588101252674& 0.0016631504549325879\\
7&-0.0010387436499409148&-0.0010286405151170262&-0.00096707869183971032\\
8& 0.0005771411154798442& 0.00058193637214214058& 0.00055921707572921596\\
9&-0.00032000941782624849&-0.00032758961661547270&-0.00032068701986538825\\
10&0.00017513356052392218&0.00018152666293358129&0.00018044775438430657\\
11&-0.000091227051122182618&-0.000095466823519254089&-0.000096030794923079510\\
12&0.000038902947970792916&0.000040963266955134492&0.000041522293589206772\\
\bottomrule
\end{tabular}}
\end{center}

The JSON candidate also contains exactly derived physical power-basis
coefficients.  They are obtained from the identity
\[
 T_{n+1}(x)=2xT_n(x)-T_{n-1}(x)
\]
using rational polynomial arithmetic, so they introduce no approximation.

\section{Type-IIc exact-dual replay jobs}
\label{app:dual-jobs}

Here \(L_m\) denotes lower-band count \(m\) and \(U_m\) upper-band count
\(m\).  Every \(T^*\) is at least the sum of the two actual caps.  The
subdivision is \(N\times N\).

\begin{longtable}{llllrr}
\toprule
Pair&Log file&\(T^*\)&\(\omega\)&\(N\)&\(n\)\\
\midrule
\endfirsthead
\toprule
Pair&Log file&\(T^*\)&\(\omega\)&\(N\)&\(n\)\\
\midrule
\endhead
\(L_3,U_5\)&\texttt{exact\_m35\_N8.log}&73/200&11/2000&8&8\\
\(L_3,U_6\)&\texttt{exact\_m36\_N8.log}&46/125&11/2000&8&9\\
\(L_4,U_4\)&\texttt{exact\_m44\_N8.log}&73/200&11/2000&8&8\\
\(L_4,U_5\)&\texttt{exact\_surgery2\_m45\_N16.log}&187/500&11/2000&16&9\\
\(L_4,U_6\)&\texttt{exact\_m46\_N8.log}&191/500&11/2000&8&10\\
\(L_5,U_4\)&\texttt{exact\_surgery3\_m54\_N16.log}&371/1000&11/2000&16&9\\
\(L_5,U_5\)&\texttt{exact\_surgery2\_m55\_N16.log}&97/250&11/2000&16&10\\
\(L_5,U_6\)&\texttt{exact\_m56\_N8.log}&99/250&11/2000&8&11\\
\(L_6,U_3\)&\texttt{exact\_m63\_N8.log}&369/1000&11/2000&8&9\\
\(L_6,U_4\)&\texttt{exact\_m64\_N8.log}&383/1000&11/2000&8&10\\
\(L_6,U_5\)&\texttt{exact\_m65\_N8.log}&397/1000&11/2000&8&11\\
\(L_6,U_6\)&\texttt{typeiic\_exact\_dual\_N32.log}&2/5&11/2000&32&12\\
\addlinespace
\(U_1,U_6\)&\texttt{exact\_u16\_N8.log}&351/1000&1/125&8&7\\
\(U_2,U_6\)&\texttt{exact\_u26\_N8.log}&351/1000&1/125&8&8\\
\(U_3,U_3\)&\texttt{exact\_u33\_N8.log}&167/500&1/125&8&6\\
\(U_3,U_4\)&\texttt{exact\_u34\_N8.log}&87/250&1/125&8&7\\
\(U_3,U_5\)&\texttt{exact\_u35\_N8.log}&181/500&1/125&8&8\\
\(U_3,U_6\)&\texttt{exact\_u36\_N8.log}&73/200&1/125&8&9\\
\(U_4,U_4\)&\texttt{exact\_u44\_N8.log}&181/500&1/125&8&8\\
\(U_4,U_5\)&\texttt{exact\_u45\_N32.log}&47/125&1/125&32&9\\
\(U_4,U_6\)&\texttt{exact\_u46\_N8.log}&379/1000&1/125&8&10\\
\(U_5,U_5\)&\texttt{exact\_surgery2\_u55\_N16.log}&19/50&1/125&16&10\\
\(U_5,U_6\)&\texttt{exact\_surgery2\_u56\_N16.log}&19/50&1/125&16&11\\
\(U_6,U_6\)&\texttt{exact\_surgery2\_u66\_N16.log}&19/50&1/125&16&12\\
\bottomrule
\end{longtable}

The three retained but no longer required mixed jobs are
\texttt{exact\_m51\_N8.log}, \texttt{exact\_m52\_N8.log}, and
\texttt{exact\_m53\_N8.log}.  The script
\texttt{analysis/verify\_final\_support.py} checks the mapping in this table,
including \(n=m+m'\), the correct \(\omega\), and \(T^*\geq B+B'\).

\begin{thebibliography}{9}

\bibitem{Maynard2015}
J. Maynard,
\emph{Small gaps between primes},
Annals of Mathematics \textbf{181} (2015), 383--413.
\href{https://doi.org/10.4007/annals.2015.181.1.7}
{doi:10.4007/annals.2015.181.1.7}.

\bibitem{Polymath2014}
D. H. J. Polymath,
\emph{Variants of the Selberg sieve, and bounded intervals containing many primes},
Research in the Mathematical Sciences \textbf{1} (2014), Article 12.
\href{https://arxiv.org/abs/1407.4897}{arXiv:1407.4897}.

\bibitem{Stadlmann2026}
J. Stadlmann,
\emph{Bounded gaps between primes},
arXiv:2608.31126v1 (31 August 2026).
\href{https://arxiv.org/abs/2608.31126}{arXiv:2608.31126}.

\end{thebibliography}

\end{document}
